\section[Cote 85 : représentations de $G_2$ dans $\mathrm{PGL}_2$, rotation d'ordre 5]{Cote 85 : représentations de $G_2$ dans $\mathrm{PGL}_2$, rotation d'ordre 5\protect\verifie{Folder85}}\label{sec:85}

La cote 85, \emph{2-polyèdres et 1-polygônes : notes manuscrites (s.d.)}, du groupe \q{Géométrie et topologie combinatoire}, étudie les polygones et polyèdres réguliers réalisés dans une forme de $\mathrm{PGL}_2$ sur une base quelconque, au moyen d'un seul invariant $A$. Pour $\mathrm{PGL}_2$, la lecture le donne par $A(v) = (\mathrm{Tr}\, v)^2/\det v - 2$ pour un relèvement $v \in \mathrm{GL}_2$ (pages~23 et~27, \lot{85}{2}{23}, \lot{85}{2}{27}). La page~30 dit qu'un élément est d'ordre~$5$ en tout point si et seulement si $\alpha^2 + \alpha - 1 = 0$, hors caractéristique~$5$ (\lot{85}{2}{30}). La page~10 énonce le cas de l'icosaèdre avec un facteur $\mathbf Z/2$ et la marge \q{laisser tomber $\mathbf Z/2$ dans l'énoncé} (\lot{85}{1}{10}). La page~40 énonce le théorème sans ce facteur : pour une forme $G$ de $\mathrm{PGL}_2$ sur un schéma $S$ et un monomorphisme $\varphi : \mathfrak A_5 \to G$, $\alpha(\varphi) = A(\varphi(\pi))$, $\pi$ un $5$-cycle, satisfait $\alpha^2 + \alpha - 1 = 0$, la catégorie des $(G, \varphi)$ est rigide, et le foncteur des classes d'isomorphisme est \q{représentable par $M = \mathrm{Spec}\,\mathbf Z[T]/(T^2 + T - 1)$} (\lot{85}{2}{40}). Le calcul qui porte ce théorème est démontré ici, sur tout anneau commutatif, ainsi qu'un modèle explicite des relations de l'icosaèdre.

Le théorème de la page~4 (\lot{85}{1}{4}) porte sur le groupe $G_2 = \langle \tau_0, \tau_1, \tau_2 \mid \tau_i^2, (\tau_0\tau_2)^2\rangle$, avec $\rho_0 = \tau_0\tau_1$ et $\rho_1 = \tau_1\tau_2$. Pour un homomorphisme $\varphi$ de $G_2$ dans une forme de $\mathrm{PGL}_2$ sur $S$ tel qu'en tout point $\rho_0^2 \neq 1$, $\rho_1^2 \neq 1$ et $\rho_0\rho_1 \neq \rho_1\rho_0$, la page énonce trois choses : 1)~les invariants $\alpha_0 = A(\varphi(\rho_0))$, $\alpha_1 = A(\varphi(\rho_1))$ sont tels que $(\alpha_0 + 2)(\alpha_1 + 2)(\alpha_0 + \alpha_1)$ est inversible ; 2)~le foncteur $\varphi \mapsto (\alpha_0, \alpha_1)$ est pleinement fidèle, c'est-à-dire que le centralisateur de l'image est trivial et que deux objets de mêmes invariants sont isomorphes ; 3)~tout couple admissible provient d'un $\varphi$. La lecture en tire que l'espace de modules est l'ouvert de $\mathbb A^2_{\mathbf Z}$ complémentaire des trois droites $\alpha_0 = -2$, $\alpha_1 = -2$, $\alpha_0 + \alpha_1 = 0$. La démonstration des pages~5 et~6 renvoie à des \q{généralités} absentes du dossier (\lot{85}{1}{5}, \lot{85}{1}{6}). Le corollaire des pages~6 et~7 prescrit les ordres de $\rho_0$ et $\rho_1$ par $\Psi_{\nu_i}(\alpha_i) = 0$, avec $\nu_i$ inversible si $\nu_i$ n'est pas premier (\lot{85}{1}{7}). Les pages~8 à~10 traitent les trois cas finis (\lot{85}{1}{8}). On démontre ici les calculs de traces sur tout anneau, la rigidité sur tout anneau, la conjugaison sur un corps qui contient une racine de $T^2 - \alpha_0 T + 1$, une forme normale explicite après adjonction d'une telle racine, et la traduction des ordres $3$, $4$, $5$ et $6$.

\noindent\emph{Conventions.} $R$ est un anneau commutatif et $g = \begin{pmatrix} a & b \\ c & d \end{pmatrix}$ une matrice de $M_2(R)$ de déterminant $\delta$ inversible ; on note $t = \mathrm{Tr}\, g$ et $A(g) = t^2\delta^{-1} - 2$. On dit que $g$ est \emph{nulle part scalaire} si $b$, $c$ et $a - d$ engendrent l'idéal unité. Un idéal est l'idéal unité si et seulement s'il n'est contenu dans aucun idéal premier ; la condition dit donc que l'image de $g$ dans $\mathrm{PGL}_2(k(s))$ est différente de $1$ en tout point $s$ de $\operatorname{Spec} R$. Comme $5$ est premier, l'image de $g$ est alors d'ordre~$5$ en tout point si et seulement si $g^5$ est scalaire.

\noindent\emph{Conventions pour $G_2$.} Un homomorphisme $G_2 \to \mathrm{PGL}_2(R)$ est donné par trois matrices $\tau_0, \tau_1, \tau_2$ de déterminants $d_0, d_1, d_2$ inversibles, telles que $\tau_0^2$, $\tau_1^2$, $\tau_2^2$ et $(\tau_0\tau_2)^2$ soient scalaires. On note $\rho_0 = \tau_0\tau_1$, $\rho_1 = \tau_1\tau_2$, $\alpha_0 = A(\rho_0)$, $\alpha_1 = A(\rho_1)$, $t_{01} = \mathrm{Tr}(\tau_0\tau_1)$, $t_{12} = \mathrm{Tr}(\tau_1\tau_2)$, $t_{02} = \mathrm{Tr}(\tau_0\tau_2)$, $T = \mathrm{Tr}(\tau_0\tau_1\tau_2)$ et $\Delta = (\alpha_0 + 2)(\alpha_1 + 2)(\alpha_0 + \alpha_1)$. Sur un corps $k$, la condition~$(*)$ de la page~4 s'écrit : $\rho_0^2$ et $\rho_1^2$ ne sont pas scalaires, et il n'existe pas de $\mu \in k$ tel que $\rho_0\rho_1 = \mu\,\rho_1\rho_0$. \q{$(*)$ en tout point} signifie que $(*)$ vaut pour les images des $\tau_i$ dans $M_2(R/\mathfrak m)$, pour tout idéal maximal $\mathfrak m$. Pour $\nu \geq 2$, on dit que $g$ est d'ordre $\nu$ en tout point si $g^\nu$ est scalaire et si $g^m$ est nulle part scalaire pour tout diviseur $m$ de $\nu$ avec $1 \leq m < \nu$.

\begin{lemme}\label{lem:85-traces}
Soient $\tau_0, \tau_1, \tau_2 \in M_2(R)$ de trace nulle et de déterminants inversibles.
\begin{enumerate}
\item $\alpha_0 + 2 = t_{01}^2 (d_0 d_1)^{-1}$ et $\alpha_1 + 2 = t_{12}^2 (d_1 d_2)^{-1}$.
\item $\rho_0\rho_1 - \rho_1\rho_0 = -T\,\tau_1$.
\item Si $t_{02} = 0$, alors $T^2 = 4 d_0 d_1 d_2 - t_{01}^2 d_2 - t_{12}^2 d_0$, donc $\alpha_0 + \alpha_1 = -T^2 (d_0 d_1 d_2)^{-1}$.
\item Si $t_{02} = 0$, $\Delta$ est inversible si et seulement si $t_{01}$, $t_{12}$ et $T$ le sont.
\end{enumerate}
\end{lemme}

\begin{proof}
(1) $\det \rho_0 = d_0 d_1$ et $A(\rho_0) = t_{01}^2 (d_0 d_1)^{-1} - 2$ par définition. (2) et (3) sont des identités polynomiales entre les coefficients, une fois $\tau_i = \begin{pmatrix} a_i & b_i \\ c_i & -a_i \end{pmatrix}$ ; pour (3), la différence des deux membres est $t_{02}$ multiplié par un polynôme explicite, et c'est ainsi que Lean la vérifie. Le membre de droite de~(3) est, au signe près, le déterminant de Gram des $\tau_i$ pour la forme $\mathrm{Tr}(XY)$, avec $t_{02} = 0$. En reportant~(3) dans $\alpha_0 + \alpha_1 = t_{01}^2(d_0d_1)^{-1} + t_{12}^2(d_1d_2)^{-1} - 4$, on obtient $-T^2(d_0d_1d_2)^{-1}$. (4) Par~(1) et~(3), $\Delta = -t_{01}^2 t_{12}^2 T^2 \cdot d_0^{-2} d_1^{-3} d_2^{-2}$, et un produit est inversible si et seulement si chaque facteur l'est.
\end{proof}

\begin{enonce}[Page 4, partie 1), les invariants\piste{85-g2-pgl2-moduli-over-z}]
Soient $\tau_0, \tau_1, \tau_2 \in M_2(R)$ de déterminants inversibles, avec $\tau_i^2$ et $(\tau_0\tau_2)^2$ scalaires.
\begin{enumerate}
\item Si $(*)$ vaut en tout point, alors $\tau_0$, $\tau_1$, $\tau_2$ et $\tau_0\tau_2$ sont de trace nulle, et $\Delta = (\alpha_0 + 2)(\alpha_1 + 2)(\alpha_0 + \alpha_1)$ est inversible.
\item Réciproquement, si $\tau_0$, $\tau_1$, $\tau_2$ et $\tau_0\tau_2$ sont de trace nulle et $\Delta$ inversible, alors $(*)$ vaut en tout point.
\end{enumerate}
\end{enonce}

\begin{proof}
(1) Fixons un idéal maximal $\mathfrak m$ et travaillons dans $k = R/\mathfrak m$. Si $\tau_0$ y était scalaire, $\rho_0^2 = c^2\tau_1^2$ serait scalaire ; de même pour $\tau_1$ (avec $\tau_0^2$) et pour $\tau_2$ (avec $\rho_1^2 = c^2\tau_1^2$). Si $\tau_0\tau_2 = c$ était scalaire, $c$ serait non nul, $\tau_2\tau_0 = c$ aussi, et $\rho_0\rho_1 = \tau_0\tau_1^2\tau_2$ et $\rho_1\rho_0 = \tau_1(\tau_2\tau_0)\tau_1$ vaudraient tous deux $c\,s$, où $\tau_1^2 = s$ ; cela contredit~$(*)$ avec $\mu = 1$. Les quatre matrices sont donc nulle part scalaires sur $R$, car un idéal qui n'est contenu dans aucun idéal maximal est l'idéal unité. Pour une matrice $g$ nulle part scalaire avec $g^2$ scalaire, Cayley--Hamilton donne $g^2 = (\mathrm{Tr}\,g)\,g - \det g$ ; les coefficients hors diagonale de $(\mathrm{Tr}\,g)\,g$ et la différence de ses coefficients diagonaux sont nuls, donc $\mathrm{Tr}\,g$ annule un idéal égal à l'idéal unité, et $\mathrm{Tr}\,g = 0$. Les quatre traces sont nulles. Revenons dans $k$. Si $t_{01} = 0$, $\rho_0$ est de trace nulle et $\rho_0^2 = -\det\rho_0$ est scalaire ; de même $t_{12} \neq 0$. Si $T = 0$, le lemme~\ref{lem:85-traces}~(2) donne $\rho_0\rho_1 = \rho_1\rho_0$. Les images de $t_{01}$, $t_{12}$ et $T$ sont donc non nulles en tout point, ces trois éléments sont inversibles, et $\Delta$ l'est par le lemme~\ref{lem:85-traces}~(4).

(2) Par le lemme~\ref{lem:85-traces}~(4), $t_{01}$, $t_{12}$ et $T$ sont inversibles, et le restent dans chaque $R/\mathfrak m$. Il suffit donc de montrer $(*)$ sur un anneau non nul où ces trois éléments sont inversibles. On a $T = \mathrm{Tr}(\rho_0\tau_2) = (\rho_{0,00} - \rho_{0,11})\,\tau_{2,00} + \rho_{0,01}\,\tau_{2,10} + \rho_{0,10}\,\tau_{2,01}$, car $\tau_2$ est de trace nulle ; $T$ appartient donc à l'idéal engendré par les coefficients hors diagonale de $\rho_0$ et la différence de ses coefficients diagonaux, et $\rho_0$ est nulle part scalaire. Si $\rho_0^2$ était scalaire, on aurait $t_{01} = \mathrm{Tr}\,\rho_0 = 0$. Même raisonnement pour $\rho_1$, avec $T = \mathrm{Tr}(\rho_1\tau_0)$. Si enfin $\rho_0\rho_1 = \mu\,\rho_1\rho_0$, multiplions à gauche par $\tau_0$ et prenons la trace : $\mathrm{Tr}(\tau_0\rho_0\rho_1) = d_0 d_1 \mathrm{Tr}\,\tau_2 = 0$, et $\mathrm{Tr}(\tau_0\rho_1\rho_0) = t_{01} T$ (identité polynomiale pour des matrices de trace nulle). Donc $\mu\,t_{01}T = 0$, $\mu = 0$, et $\rho_0\rho_1 = 0$, ce qui contredit $\det(\rho_0\rho_1) = d_0 d_1^2 d_2$ inversible.
\end{proof}

\begin{remarque}
Chaque facteur de $\Delta$ répond à une condition de~$(*)$ : $\alpha_0 + 2$ inversible équivaut à $t_{01}$ inversible, c'est-à-dire $\rho_0^2 \neq 1$ ; $\alpha_1 + 2$ à $\rho_1^2 \neq 1$ ; et $\alpha_0 + \alpha_1$, au vu de $\rho_0\rho_1 - \rho_1\rho_0 = -T\tau_1$, à la non-commutation. L'énoncé vaut sur tout anneau, sans hypothèse sur $2$ : en caractéristique~$2$ les involutions sont unipotentes, et $\tau_0$, $\tau_2$ commutent au lieu d'anticommuter, sans que les calculs changent. L'hypothèse \q{$(\tau_0\tau_2)^2$ scalaire} sert à obtenir $t_{02} = 0$ ; sans elle, l'identité~(3) du lemme a un terme en $t_{02}$.
\end{remarque}

\begin{lemme}\label{lem:85-commutant}
Si $X \in M_2(R)$ est nulle part scalaire et $gX = Xg$, il existe $p, q \in R$ avec $g = p + qX$.
\end{lemme}

\begin{proof}
Écrivons $X = \begin{pmatrix} a & b \\ c & d \end{pmatrix}$, $g = \begin{pmatrix} e & f \\ h & k \end{pmatrix}$. La commutation donne $fc = bh$, $b(e - k) = f(a - d)$ et $h(a - d) = c(e - k)$ : les vecteurs $(f, h, e - k)$ et $(b, c, a - d)$ sont proportionnels au sens où leurs mineurs d'ordre~$2$ sont nuls. Soit $r_1 b + r_2 c + r_3(a - d) = 1$. Posons $q = r_1 f + r_2 h + r_3(e - k)$ et $p = e - qa$. Alors $f = f(r_1 b + r_2 c + r_3(a-d)) = qb$ par les relations, de même $h = qc$ et $e - k = q(a - d)$, d'où $g = p + qX$.
\end{proof}

\begin{enonce}[Page 4, partie 2), rigidité]
Soient $\tau_0, \tau_1, \tau_2 \in M_2(R)$ de déterminants inversibles, de trace nulle, avec $t_{02} = 0$ et $\Delta$ inversible. Si $g \in \mathrm{GL}_2(R)$ et $\mu_0, \mu_1, \mu_2 \in R$ vérifient $g\tau_i = \mu_i\,\tau_i g$ pour $i = 0, 1, 2$, alors $g$ est scalaire.
\end{enonce}

\begin{proof}
On a $g\rho_0 = \mu_0\mu_1\,\rho_0 g$, d'où $t_{01} = \mathrm{Tr}(g^{-1}g\rho_0) = \mu_0\mu_1\,\mathrm{Tr}(g^{-1}\rho_0 g) = \mu_0\mu_1 t_{01}$, et $\mu_0\mu_1 = 1$ puisque $t_{01}$ est inversible. De même $\mu_1\mu_2 = 1$ : $g$ commute à $\rho_0$ et à $\rho_1$. Comme $T$ est inversible, $\rho_0$ est nulle part scalaire (démonstration de l'énoncé précédent), et le lemme~\ref{lem:85-commutant} donne $g = p + q\rho_0$. La commutation avec $\rho_1$ donne $q(\rho_0\rho_1 - \rho_1\rho_0) = 0$, soit $qT\tau_1 = 0$ ; en multipliant par $\tau_1$, $qTd_1 = 0$, donc $q = 0$ et $g = p$.
\end{proof}

\noindent\emph{La forme normale.} Pour $\ell, b \in R$, posons
\[
\rho_0^\circ = \begin{pmatrix} \ell & 1 \\ 0 & 1 \end{pmatrix}, \quad
\rho_1^\circ = \begin{pmatrix} 0 & 1 \\ -b & b \end{pmatrix}, \quad
\tau_1^\circ = \begin{pmatrix} b & 1 - \ell - b \\ b(1 - \ell) & -b \end{pmatrix}, \quad
\tau_0^\circ = \rho_0^\circ\tau_1^\circ, \quad \tau_2^\circ = \tau_1^\circ\rho_1^\circ .
\]
La matrice $\rho_0^\circ$ a pour trace $\ell + 1$ et pour déterminant $\ell$, et $A(\rho_0^\circ) = \ell + \ell^{-1}$ ; $\rho_1^\circ$ a pour trace et pour déterminant $b$, et $A(\rho_1^\circ) = b - 2$. La matrice $\tau_1^\circ$ est, à un scalaire près, l'unique matrice de trace nulle orthogonale à $\rho_0^\circ$ et à $\rho_1^\circ$ pour la forme $\mathrm{Tr}(XY)$.

\begin{enonce}[Page 4, partie 3), une forme normale]
Soient $\alpha_0, \alpha_1 \in R$ avec $\Delta = (\alpha_0 + 2)(\alpha_1 + 2)(\alpha_0 + \alpha_1)$ inversible, $\ell \in R$ une racine de $T^2 - \alpha_0 T + 1$ et $b = \alpha_1 + 2$. Alors $\tau_0^\circ$, $\tau_1^\circ$, $\tau_2^\circ$ sont de déterminants inversibles, de trace nulle, avec $\mathrm{Tr}(\tau_0^\circ\tau_2^\circ) = 0$ (de sorte que $(\tau_i^\circ)^2$ et $(\tau_0^\circ\tau_2^\circ)^2$ sont scalaires), $A(\tau_0^\circ\tau_1^\circ) = \alpha_0$, $A(\tau_1^\circ\tau_2^\circ) = \alpha_1$, et $(*)$ vaut en tout point. Pour tout anneau $R$ et tout tel couple, cela s'applique sur $R[\ell] = R[T]/(T^2 - \alpha_0 T + 1)$, libre de rang~$2$ sur $R$.
\end{enonce}

\begin{proof}
On a $\ell(\alpha_0 - \ell) = 1$, donc $\ell$ est inversible. Les traces de $\tau_1^\circ$, $\tau_0^\circ$, $\tau_2^\circ$ et $\tau_0^\circ\tau_2^\circ$ sont nulles par calcul direct. Ensuite $\det\tau_1^\circ = -b\big((1-\ell)^2 + \ell b\big)$ et, comme $\ell^2 + 1 = \alpha_0\ell$, $(1 - \ell)^2 + \ell b = \ell(\alpha_0 + \alpha_1)$ ; ce déterminant est inversible, et $\det\tau_0^\circ = \ell\det\tau_1^\circ$, $\det\tau_2^\circ = b\det\tau_1^\circ$. Comme $\tau_1^\circ$ est de trace nulle, $(\tau_1^\circ)^2 = s$ avec $s = -\det\tau_1^\circ$ inversible, d'où $\tau_0^\circ\tau_1^\circ = s\rho_0^\circ$ et $\tau_1^\circ\tau_2^\circ = s\rho_1^\circ$. L'invariant ne dépend pas du relèvement : $A(\tau_0^\circ\tau_1^\circ) = (\ell + 1)^2\ell^{-1} - 2 = \alpha_0$, car $(\ell + 1)^2 = (\alpha_0 + 2)\ell$, et $A(\tau_1^\circ\tau_2^\circ) = b^2 b^{-1} - 2 = \alpha_1$. Enfin $\Delta$ est inversible, et la partie~(2) de l'énoncé des invariants donne $(*)$ en tout point. Dans $R[T]/(T^2 - \alpha_0 T + 1)$, la classe de $T$ est racine par construction.
\end{proof}

\begin{remarque}
La forme normale demande une racine $\ell$ de $T^2 - \alpha_0 T + 1$. Sur $R$ lui-même, une réalisation dans le $\mathrm{PGL}_2$ déployé n'existe pas toujours : pour $R = \mathbf R$ et $(\alpha_0, \alpha_1) = (-1, c_5)$, elle donnerait un $\mathfrak A_5$ dans $\mathrm{PGL}_2(\mathbf R)$, qui n'en contient pas (fait classique, non formalisé ici). C'est pourquoi la page~4 parle de formes de $\mathrm{PGL}_2$. Le passage de $R[\ell]$ à une forme sur $R$, par descente le long de $R \to R[\ell]$ en s'appuyant sur la rigidité, n'est pas formalisé.
\end{remarque}

\begin{enonce}[Page 4, partie 2), conjugaison sur un corps]
Soient $K$ un corps et $\tau_0, \tau_1, \tau_2 \in \mathrm{GL}_2(K)$ de trace nulle, avec $t_{02} = 0$ et $\Delta \neq 0$. Soit $\ell \in K$ une racine de $T^2 - \alpha_0 T + 1$ et $b = \alpha_1 + 2$. Il existe $g \in \mathrm{GL}_2(K)$ et $c_0, c_1, c_2 \in K$ tels que $\tau_i g = c_i\, g\tau_i^\circ$ pour $i = 0, 1, 2$. En particulier, deux tels triplets de mêmes invariants $(\alpha_0, \alpha_1)$ sont conjugués dans $\mathrm{PGL}_2(K)$.
\end{enonce}

\begin{proof}
Le lemme~\ref{lem:85-traces} donne $t_{01}$, $t_{12}$, $T$ non nuls. Posons $X = (\ell + 1)t_{01}^{-1}\rho_0$ et $Y = b\,t_{12}^{-1}\rho_1$. On a $\ell \neq 0$, $\ell + 1 \neq 0$ car $(\ell+1)^2 = (\alpha_0 + 2)\ell$, et $b \neq 0$. Alors $\mathrm{Tr}\,X = \ell + 1$, $\det X = (\ell+1)^2 t_{01}^{-2} d_0 d_1 = \ell$, $\mathrm{Tr}\,Y = \det Y = b$, $\mathrm{Tr}(XY)$ est un multiple de $\mathrm{Tr}(\tau_0\tau_1^2\tau_2) = -d_1 t_{02} = 0$, et $\det(XY - YX)$ est un multiple non nul de $\det(T\tau_1) = T^2 d_1$.

Comme $\det(X - \ell) = \ell^2 - \ell(\ell + 1) + \ell = 0$, il existe $v \neq 0$ avec $Xv = \ell v$. Posons $w = -b^{-1}Yv$ et soit $g$ la matrice de colonnes $v$, $w$. Par définition $Yv = -bw$ ; par Cayley--Hamilton, $Y^2 = bY - b$, d'où $Yw = v + bw$. L'identité $XY + YX = (\mathrm{Tr}\,X)Y + (\mathrm{Tr}\,Y)X + \mathrm{Tr}(XY) - \mathrm{Tr}\,X\,\mathrm{Tr}\,Y$ donne $XYv = Yv - bv$, soit $Xw = v + w$. Ces quatre égalités disent $Xg = g\rho_0^\circ$ et $Yg = g\rho_1^\circ$. Si $\det g = 0$, $w$ est proportionnel à $v$, qui est alors vecteur propre commun de $X$ et $Y$ ; $(XY - YX)v = 0$ et $\det(XY - YX) = 0$, ce qui est exclu. En Lean, cette étape passe par les identités $b\,v_i\,\big((XY - YX)v\big)_j \in (\det g)$ modulo les relations, et le choix d'une coordonnée $v_i \neq 0$.

Soit $M = g^{-1}\tau_1 g$. Sa trace est nulle ; $\mathrm{Tr}(M\rho_0^\circ) = \mathrm{Tr}(\tau_1 X)$ est un multiple de $\mathrm{Tr}(\tau_1\tau_0\tau_1) = -d_1\mathrm{Tr}\,\tau_0 = 0$, et de même $\mathrm{Tr}(M\rho_1^\circ) = 0$. Ces trois équations linéaires imposent $M = M_{00}b^{-1}\,\tau_1^\circ$. Enfin $\tau_0 = (-d_1)^{-1}(\ell+1)^{-1}t_{01}\,X\tau_1$ et $\tau_2 = (-d_1)^{-1}b^{-1}t_{12}\,\tau_1 Y$, d'où $\tau_0 g \in K\,g\rho_0^\circ\tau_1^\circ = K\,g\tau_0^\circ$ et $\tau_2 g \in K\,g\tau_1^\circ\rho_1^\circ = K\,g\tau_2^\circ$. Pour deux triplets $\tau$, $\sigma$ de mêmes invariants, avec $g$, $g'$ ainsi obtenus, $h = g g'^{-1}$ vérifie $\tau_i h = c_i c_i'^{-1}\,h\sigma_i$.
\end{proof}

\begin{remarque}
La conjugaison est démontrée sur un corps qui contient $\ell$, par exemple algébriquement clos. Sur un corps quelconque, la rigidité et un argument de descente galoisienne la donneraient aussi ; ce passage n'est pas formalisé, ni la conjugaison sur un anneau local ou une base quelconque. La matrice $\rho_0^\circ$ couvre d'un seul coup le cas semi-simple ($\ell \neq 1$) et le cas unipotent ($\ell = 1$, $\alpha_0 = 2$, possible hors caractéristique~$2$).
\end{remarque}

\begin{enonce}[Pages 6 et 7, 29 à 31, les ordres]
Soit $g \in M_2(R)$ nulle part scalaire, de trace $t$ et de déterminant $d$ inversible, et $\alpha = A(g)$. Définissons $x_0 = 0$, $x_1 = 1$, $x_{n+2} = t x_{n+1} - d x_n$, et $P_0 = 0$, $P_1 = 1$, $P_{n+2} = P_{n+1} - P_n$ pour $n$ pair, $P_{n+2} = (\alpha + 2)P_{n+1} - P_n$ pour $n$ impair.
\begin{enumerate}
\item $g^{n+1} = x_{n+1}\,g - d\,x_n$ ; $g^{n+1}$ est scalaire si et seulement si $x_{n+1} = 0$, et nulle part scalaire si et seulement si $x_{n+1}$ est inversible.
\item $x_{2k+1} = d^k P_{2k+1}(\alpha)$ et $x_{2k+2} = t\,d^k P_{2k+2}(\alpha)$. Donc $g^{2k+1}$ est scalaire si et seulement si $P_{2k+1}(\alpha) = 0$, et, si $t$ est inversible, $g^{2k+2}$ est scalaire si et seulement si $P_{2k+2}(\alpha) = 0$.
\item $P_3 = \alpha + 1$, $P_4 = \alpha$, $P_5 = \alpha^2 + \alpha - 1$, $P_6 = \alpha^2 - 1$, $P_7 = \alpha^3 + \alpha^2 - 2\alpha - 1$.
\item $g$ est d'ordre~$3$ en tout point si et seulement si $\alpha = -1$ ; d'ordre~$4$ en tout point si et seulement si $\alpha = 0$ et $2$ est inversible ; d'ordre~$6$ en tout point si et seulement si $\alpha = 1$ et $6$ est inversible.
\end{enumerate}
\end{enonce}

\begin{proof}
(1) Par récurrence, avec $g^2 = tg - d$ : $g^{n+2} = (x_{n+1}g - dx_n)g = (tx_{n+1} - dx_n)g - dx_{n+1}$. Pour $g$ nulle part scalaire, $xg + y$ est scalaire si et seulement si $x = 0$ (argument d'idéal), et nulle part scalaire si et seulement si $x$ est inversible : ses coefficients hors diagonale et la différence diagonale sont $x$ fois ceux de $g$, et un idéal maximal contient $xb$, $xc$, $x(a-d)$ si et seulement s'il contient $x$, ou $b$, $c$ et $a - d$. (2) Par récurrence sur $k$, avec $t^2 = (\alpha + 2)d$ : $x_{2k+3} = t\cdot t d^k P_{2k+2} - d\cdot d^k P_{2k+1} = d^{k+1}\big((\alpha+2)P_{2k+2} - P_{2k+1}\big)$ et $x_{2k+4} = t d^{k+1}P_{2k+3} - d\cdot t d^k P_{2k+2} = t d^{k+1}(P_{2k+3} - P_{2k+2})$. (3) Calcul direct. (4) Ordre~$3$ : $g^3$ scalaire équivaut à $P_3(\alpha) = 0$. Ordre~$4$ : $g^2$ nulle part scalaire équivaut à $x_2 = t$ inversible, et $g^4$ scalaire à $t\,d\,\alpha = 0$, soit $\alpha = 0$ ; alors $t^2 = 2d$, donc $2$ est inversible ; réciproquement $\alpha = 0$ et $2$ inversible donnent $t^2 = 2d$ inversible. Ordre~$6$ : il faut $t$ inversible, $x_3 = d(\alpha + 1)$ inversible et $t d^2(\alpha - 1)(\alpha + 1) = 0$, soit $\alpha = 1$ ; alors $2 = \alpha + 1$ est inversible et $t^2 = 3d$ donne $3$ inversible ; la réciproque est le même calcul.
\end{proof}

\begin{remarque}
La suite $x_n$ est celle du lemme~\ref{lem:85-puissance}, et $P_5 = \Psi_5$ redonne le calcul de l'ordre~$5$. Les valeurs de~(3) sont $\Psi_3$, $\Psi_4$, $\Psi_5$, $\Psi_3\Psi_6$ et $\Psi_7$ : ce sont les $F_n$ de la page~29 pour $n \leq 7$. L'identification $P_n = F_n = \prod_{d \mid n,\, d \geq 3}\Psi_d$ pour tout $n$ n'est pas formalisée. Le point~(4) confirme la condition~c) du corollaire des pages~6 et~7 pour $\nu = 4$ et $\nu = 6$ : dans $\mathrm{PGL}_2$, l'ordre~$6$ demande $6$ inversible, alors que, dans $\mathrm{SL}_2$, le cas~B) de la page~31 ne demande que $2$ ; un élément de $\mathrm{PGL}_2(k)$ d'invariant $1$ en caractéristique~$3$ a une trace nulle et est une involution. Pour les cas finis des pages~8 à~10, $\Delta(-1, -1) = -2$, $\Delta(0, -1) = -2$ et, si $\alpha^2 + \alpha - 1 = 0$, $\Delta(-1, \alpha) = -1$ : les deux premiers demandent $2$ inversible, le troisième est inversible sur toute base, comme le dit la lecture.
\end{remarque}

\begin{enonce}[Pages 10, 30 et 40, l'invariant d'une rotation d'ordre~$5$\piste{85-a5-pgl2-moduli-spec-z-phi}]
Soit $g \in M_2(R)$ de déterminant inversible.
\begin{enumerate}
\item $A(\mu g) = A(g)$ pour tout $\mu \in R^\times$ : $A$ ne dépend que de l'image de $g$ dans $\mathrm{PGL}_2(R)$.
\item Si $g$ est nulle part scalaire et $g^5$ scalaire, alors $A(g)^2 + A(g) - 1 = 0$.
\item Réciproquement, si $A(g)^2 + A(g) - 1 = 0$, alors $g^5$ est scalaire.
\item $A(g^2) = A(g)^2 - 2$ ; en particulier, si $A(g)^2 + A(g) - 1 = 0$, alors $A(g^2) = -1 - A(g)$ (corollaire de la page~40).
\end{enumerate}
\end{enonce}

\begin{lemme}\label{lem:85-puissance}
Pour toute matrice $g \in M_2(R)$, $g^5 = x_5\, g + y_5$, avec $x_5 = t^4 - 3\delta t^2 + \delta^2$ et $y_5 = -\delta(t^3 - 2\delta t)$. Si $\delta$ est inversible, $x_5 = \delta^2\,(A(g)^2 + A(g) - 1)$.
\end{lemme}

\begin{proof}
Par Cayley--Hamilton, $g^2 = t g - \delta$. Si $g^n = x_n g + y_n$, alors $g^{n+1} = (t x_n + y_n) g - \delta x_n$. Partant de $x_1 = 1$, $y_1 = 0$, on trouve $x_2 = t$, $x_3 = t^2 - \delta$, $x_4 = t^3 - 2\delta t$, $y_4 = -\delta(t^2 - \delta)$, puis $x_5 = t^4 - 3\delta t^2 + \delta^2$ et $y_5 = -\delta x_4$. La vérification en Lean se fait coefficient par coefficient. Posons $u = t^2\delta^{-1}$, de sorte que $A(g) = u - 2$. Alors $A(g)^2 + A(g) - 1 = u^2 - 3u + 1$, et $\delta^2(u^2 - 3u + 1) = t^4 - 3\delta t^2 + \delta^2$.
\end{proof}

\begin{proof}[Démonstration de l'énoncé]
(1) $\mathrm{Tr}(\mu g) = \mu t$ et $\det(\mu g) = \mu^2 \delta$, donc $(\mu t)^2 (\mu^2\delta)^{-1} = t^2 \delta^{-1}$.

(2) Écrivons $g^5 = \gamma \cdot 1$. Par le lemme~\ref{lem:85-puissance}, $x_5 g = (\gamma - y_5)\cdot 1$. En lisant les coefficients hors diagonale et la différence des coefficients diagonaux, $x_5 b = x_5 c = x_5 (a - d) = 0$. L'ensemble des $r \in R$ tels que $x_5 r = 0$ est un idéal ; il contient $b$, $c$, $a - d$, donc l'idéal unité, et $x_5 = 0$. Comme $\delta^2$ est inversible, le lemme donne $A(g)^2 + A(g) - 1 = 0$.

(3) Si $A(g)^2 + A(g) - 1 = 0$, le lemme donne $x_5 = 0$ et $g^5 = y_5$.

(4) $\mathrm{Tr}(g^2) = t^2 - 2\delta$ et $\det(g^2) = \delta^2$, d'où $A(g^2) = (t^2 - 2\delta)^2\delta^{-2} - 2 = (u - 2)^2 - 2 = A(g)^2 - 2$. Si $A(g)^2 = 1 - A(g)$, cela vaut $-1 - A(g)$.
\end{proof}

\begin{enonce}[Page 40, un modèle des relations de l'icosaèdre]
Soit $\zeta \in R$ tel que $\zeta^4 + \zeta^3 + \zeta^2 + \zeta + 1 = 0$, et posons
\[
  s = \begin{pmatrix} 0 & -1 \\ 1 & 0 \end{pmatrix}, \qquad p = \begin{pmatrix} \zeta & 1 \\ 0 & \zeta^4 \end{pmatrix}.
\]
Alors $s^2 = -1$, $(sp)^3 = -1$ et $p^5$ est scalaire ; $s$, $p$ et $sp$ sont nulle part scalaires ; $\det p = 1$ et $A(p) = \zeta^2 + \zeta^3$ est racine de $T^2 + T - 1$. C'est le cas en particulier pour $R = \mathbf Z[\zeta_5] = \mathbf Z[X]/(X^4 + X^3 + X^2 + X + 1)$ et $\zeta$ la classe de $X$.
\end{enonce}

\begin{proof}
On a $\zeta^5 - 1 = (\zeta - 1)(\zeta^4 + \zeta^3 + \zeta^2 + \zeta + 1) = 0$, donc $\zeta$ est inversible d'inverse $\zeta^4$, et $\det p = \zeta^5 = 1$. Le calcul direct donne $s^2 = -1$. Le produit $sp = \begin{pmatrix} 0 & -\zeta^4 \\ \zeta & 1 \end{pmatrix}$ a pour trace $1$ et pour déterminant $\zeta^5 = 1$, donc $(sp)^2 = sp - 1$, puis $(sp)^3 = sp(sp - 1) = (sp - 1) - sp = -1$. Ensuite $\mathrm{Tr}\, p = \zeta + \zeta^4$ et $A(p) = (\zeta + \zeta^4)^2 - 2 = \zeta^2 + 2\zeta^5 + \zeta^8 - 2 = \zeta^2 + \zeta^3$. On a $(\zeta^2 + \zeta^3)^2 + \zeta^2 + \zeta^3 - 1 = \zeta^4 + 2\zeta^5 + \zeta^6 + \zeta^2 + \zeta^3 - 1 = \zeta^4 + \zeta^3 + \zeta^2 + \zeta + 1 = 0$, et le point~(3) de l'énoncé précédent donne $p^5$ scalaire. Le coefficient $(0,1)$ vaut $-1$ pour $s$, $1$ pour $p$ et $-\zeta^4$ pour $sp$ ; c'est chaque fois un inversible, et les trois matrices sont nulle part scalaires. Dans $\mathbf Z[X]/(X^4 + X^3 + X^2 + X + 1)$, la classe de $X$ est racine du polynôme par construction.
\end{proof}

\begin{remarque}
Le point~(2) vaut sur tout anneau commutatif, sans hypothèse sur $2$, $3$ ou $5$, ni de réduction. L'exclusion de la caractéristique~$5$ à la page~30 concerne le sens réciproque, l'ordre exact $5$ : en caractéristique~$5$, $T^2 + T - 1 = (T - 2)^2$, et $A(g) = 2$ vaut aussi pour $g = 1$. Avec l'hypothèse \q{nulle part scalaire} et $5$ premier, les points~(2) et~(3) disent qu'un élément nulle part scalaire de déterminant inversible est d'ordre~$5$ en tout point si et seulement si $A(g)^2 + A(g) - 1 = 0$, en toute caractéristique. Dans le modèle, en caractéristique~$5$, $\zeta = 1$ convient, $p$ est unipotent et $A(p) = 2$ : c'est la rotation unipotente que la lecture signale pour $\mathfrak A_5 \simeq \mathrm{PSL}_2(\mathbf F_5)$.
\end{remarque}

\begin{remarque}
Les relations $s^2$, $(sp)^3$, $p^5$ scalaires sont celles du groupe triangulaire de type $(2, 3, 5)$, qui est $\mathfrak A_5$. En tout point, le sous-groupe de $\mathrm{PGL}_2(k(s))$ engendré par les images de $s$ et $p$ est un quotient de $\mathfrak A_5$ différent de $1$, donc isomorphe à $\mathfrak A_5$, ce groupe étant simple. Ce passage n'est pas formalisé. Le modèle est dans le $\mathrm{PGL}_2$ déployé sur $\mathbf Z[\zeta_5]$, et non sur $\mathbf Z[T]/(T^2 + T - 1)$. Sur $\mathbf Q(\sqrt 5)$, corps réel où $-1$ n'est pas somme de deux carrés, le $\mathrm{PGL}_2$ déployé ne contient pas $\mathfrak A_5$ d'après le critère que la lecture cite de mémoire (note des pages~36 et~37) ; un modèle sur $\mathbf Z[\varphi]$ demande donc une forme non déployée, et c'est pourquoi le théorème de la page~40 porte sur les formes de $\mathrm{PGL}_2$.
\end{remarque}

Ne sont pas démontrés ici : la rigidité de la catégorie des $(G, \varphi)$ et la représentabilité de son foncteur des classes d'isomorphisme par $\operatorname{Spec} \mathbf Z[T]/(T^2 + T - 1)$ (page~40), ni l'existence d'un plongement de $\mathfrak A_5$ dans une forme de $\mathrm{PGL}_2$ sur cette base ; pour le théorème de la page~4, le passage des formes de $\mathrm{PGL}_2$ aux matrices de $\mathrm{GL}_2$, qui n'a lieu que localement, la descente de la forme normale de $R[\ell]$ à une forme sur $R$, la conjugaison sur un corps qui ne contient pas $\ell$ et sur une base quelconque, et donc la représentabilité par l'ouvert de $\mathbb A^2_{\mathbf Z}$ ; l'identification $P_n = F_n$ pour tout $n$ et le théorème de la page~31 pour $n$ quelconque ; l'équivalence entre \q{nulle part scalaire} et \q{différent de $1$ en tout point} pour une forme quelconque ; l'homomorphisme $\mathfrak A_5 \to \mathrm{PGL}_2(\mathbf Z[\zeta_5])$ défini par le modèle et son injectivité sur les fibres ; les autres énoncés de la cote.
